\documentclass[11pt,a4paper]{article}
\usepackage[margin=25mm]{geometry}
\usepackage{amsmath,amssymb,amsthm}
\setlength{\emergencystretch}{2em}
\newtheorem{lemma}{Lemma}
\newtheorem{theorem}[lemma]{Theorem}
\newcommand{\gen}[1]{\langle#1\rangle}
\newcommand{\Alt}{\operatorname{Alt}}
\title{An interior cut cannot first share the exceptional crossing}
\author{Human Compression project: L044 continuation}
\date{}
\begin{document}\maketitle
Retain the full critical six-point-head hypotheses and their paid
predecessor proofs. In particular $G\leq\Alt(\Omega)$ is transitive,
$|\Omega|=n=2r\geq18$, its distinguished involutions form a connected
string C-marking, and all maximal omissions are intransitive.
The actual first three generators form the supported six-point head
on $D$. The critical tail is the actual natural alternating group
on $(\Omega\setminus D)\cup\{6\}$, with the notation of
\emph{Pure alternating kernels in the critical prefixes}.
Write $P_k=G_{\leq k}$ and let $C_k$ be its orbit containing $D$.
We retain the entire proof of that note, including
\[
 \Alt(C_k)\leq P_k,\qquad
 P_k=\{(s,u)\in\operatorname{Sym}(C_k)\times J_k:
                         \operatorname{sgn}s=\operatorname{sgn}u\},
                                                        \tag{1}
\]
where $J_k$ is the literal image on $\Omega\setminus C_k$.
All eligible edges below refer to full ambient omissions.

\begin{lemma}[The two sides of a later unique cut]\label{sides}
Suppose $j\geq5$ has exactly one eligible edge. Its two omission
orbits are $C=C_{j-1}$ and $A=\Omega\setminus C$.
Every generator of index greater than $j$ fixes $C$ pointwise.
Consequently every shared label at this cut is smaller than $j$;
if there is a first shared label, it is at least $3$.
\end{lemma}
\begin{proof}
Put $P=G_{<j}$ and $Q=G_{>j}$. These groups commute.
By (1), $P$ contains the actual natural $\Alt(C_{j-1})$,
on at least eight points. Every element of $Q$ therefore preserves
that unique nontrivial orbit and fixes it pointwise: a permutation
centralizing a natural alternating group of degree at least four
is trivial on its support. Thus the $G_j=PQ$ orbit containing
$D$ is exactly $C_{j-1}$. A unique eligible edge and transitivity
of $G$ imply that $G_j$ has exactly two orbits, since adding that
one cross-orbit edge must connect all of them. This proves the
orbit assertion. The head fixes $\Omega\setminus D$, so its three
labels cannot be shared across this cut. The fixation by $Q$
excludes every label greater than $j$ as well.
\end{proof}

\begin{theorem}[Exclusion of the interior first-shared case]\label{exclude}
Suppose $5\leq j<r-1$ has exactly one eligible edge.
Then its first shared label cannot be $3$.
\end{theorem}
\begin{proof}
Assume the first shared label is $3$. Reverse the entire string,
writing $s_t=\rho_{r-1-t}$, and put
\[
 i=r-1-j\geq1,\qquad h=r-4,\qquad d=h-i=j-3\geq2.
\]
In this reversed order the unique crossing has index $i$, and
the last shared label is $h$. Its earlier subgroup
$L=\gen{s_t:t<i}=G_{>j}$ fixes $C$ pointwise by Lemma~\ref{sides}.
The tail past $h$ is precisely the original supported head on $D$.
It fixes $A$, and the size of its fixed set is therefore exactly
\[
                       F=n-6=2h+2.                     \tag{2}
\]

If $L$ is transitive on $A$, the complete transitive-side support
lemma gives $h=i+1$, contrary to $d\geq2$. We may consequently
use the complete intransitive-side orbit and support lemmas.
Their hypotheses hold literally: $i$ is interior; the second
orbit $C$ contains $D$ and is nontrivial; $L$ fixes $C$;
and the tail past the last shared label fixes $A$.
Let $k$ be the common size of the faithful $L$-orbits in $A$.
Those lemmas give $k\geq i+1$ and the following bounds, with
their full proofs retained.

For $d=2$, if $s_h$ exchanges the first active $L$-orbit with
another one, then
\[
 F\geq3k+2\geq3i+5.
\]
Together with (2), this forces $i\leq1$, and hence
$r=h+4=i+6\leq7$, contrary to $r\geq9$.
If that active orbit is preserved nontrivially instead, the
complete internal-orbit lemma gives $F\geq2h+3$, directly
contradicting (2). These are the two cases of the paid
distance-to-an-active-orbit argument when $d=2$.

For $d\geq3$, that argument and the fixed-prefix path bound give
\[
 F\geq |A|+d\geq dk+d\geq d(i+2).
\]
Equation (2) becomes $i(d-2)\leq2$. Since $i\geq1$ and
$r=i+d+4\geq9$, its only possibilities are
\[
 r=9,\quad (i,d)=(2,3)\quad\hbox{or}\quad(i,d)=(1,4).   \tag{3}
\]
Both have $F=12$, and all the preceding inequalities are
equalities. In the first case $k=3$ and $A$ consists of three
$L$-orbits of size three. The two independent involutions
generating $L$ act faithfully on each orbit; each restriction
is thus a transposition. Since $L$ fixes $C$, each generator
is a product of three transpositions on $\Omega$, contradicting
$G\leq\Alt(\Omega)$.

It remains to exclude $(i,d)=(1,4)$. Here $j=7$, $|A|=8$,
and $L=\gen{\rho_8}$ has four two-point orbits in $A$.
Write $\lambda=\rho_8$: it flips each of these four pairs and
fixes $C$. The complete orbit-distance proof starts at the pair
containing the crossing endpoint and reaches a pair on which
$s_h=\rho_3$ is nontrivial, using at least $d-1=3$ intervening
labels. There are exactly four pairs. A shortest such path
therefore visits all four, and $\rho_3$ fixes the first three
pointwise. It cannot exchange the last pair with another pair,
because those earlier pairs are fixed. It follows that the
restriction $\rho_3|_A$ is exactly the transposition on the last
pair. This conclusion uses the full active-orbit argument,
not merely its numerical inequality.

Now take $P=G_{<7}=P_6$ and let $J$ be its image on $A$.
Every element of $P$ commutes with $\lambda$, so $J$ preserves
the four pairs. The group $G_7=P\gen\lambda$ is transitive on
$A$, whereas $\lambda$ fixes each pair setwise. Hence $J$ is
transitive on the set of four pairs. It contains the single
pair flip $\rho_3|_A$. Conjugating that flip by $J$ supplies the
flip on every pair; their product is $\lambda|_A$. This product
is even, being a product of four disjoint transpositions.
Formula (1), with $C=C_6$ and $J=J_6$, now puts the literal
permutation $(1_C,\lambda|_A)=\rho_8$ in $P_6$.
But actual ambient IP gives
$P_6\cap\gen{\rho_8}=1$, since the distinguished index sets are
disjoint. This is the final contradiction in (3), proving
the theorem.
\end{proof}

\paragraph{Attribution and exact remaining scope.}
The source route is the exceptional exit from the reversed
transitive-tail argument in Cameron--Fernandes--Leemans--Mixer,
\emph{Highest rank of a polytope for $A_n$}, arXiv:1605.09173v1,
Propositions 5.26 and 5.28, physical pages 34 and 36--37.
The complete transitive-side, intransitive-side, active-orbit,
internal-orbit and fixed-prefix proofs, with every classical
and finite input, remain paid. The near-equality count and
the final extraction of the even four-pair flip are the new
argument here. The prefix-kernel proof attributes the initial
parity-fibre suggestion to the root collaborator.

The endpoint $j=r-1$ is excluded from Theorem~\ref{exclude};
it has $i=0$ after reversal, so the interior orbit bounds do
not apply. Other first shared indices, the opposite heads,
the terminal singleton and complete L044 recovery remain
separate. No computational, formal, whole-original, conjecture
or numerical-compression credit is claimed.
\end{document}
